\documentclass[11pt]{amsart}
\usepackage{calc,amssymb,amsmath,ulem,stmaryrd,fullpage}
\usepackage{enumerate}
\usepackage{alltt}
\RequirePackage[dvipsnames,usenames]{color}
\let\mffont=\sf
\normalem
%\input{kmacros3.sty}
\input{xy}
\xyoption{all}
\usepackage{hyperref}
\usepackage{graphicx}
\usepackage[all,cmtip]{xy}
\usepackage{verbatim}
\usepackage[left=0.95in,top=0.95in,right=0.95in,bottom=0.95in]{geometry}
\newcommand{\tauCohomology}{T}
\newcommand{\FCohomology}{S}


\theoremstyle{theorem}
%\newtheorem*{mainthm*}{Main Theorem}

\newcommand{\note}[1]{\marginpar{\sffamily\tiny #1}}

\def\todo#1{\textcolor{Mahogany}%
{\sffamily\footnotesize\newline{\color{Mahogany}\fbox{\parbox{\textwidth-15pt}{#1}}}\newline}}

\renewcommand{\O}{\mathcal O}
\newcommand{\ie}{{\itshape i.e.} }
\newcommand{\roundup}[1]{\lceil #1 \rceil}
\newcommand{\rounddown}[1]{\lfloor #1 \rfloor}
\renewcommand{\to}[1][]{\xrightarrow{\ #1\ }}
\newcommand{\onto}[1][]{\protect{\xrightarrow{\ #1\ }\hspace{-0.8em}\rightarrow}}
\newcommand{\into}[1][]{\lhook \joinrel \xrightarrow{\ #1\ }}
%\newcommand{DBDual}[1]{{\underline{\omega}_{#1}}}
\newcommand{\DBDual}[1]{{{\uuline \omega}{}^{\mydot}_{#1}}}
\newcommand{\Tt}{{\mathfrak{T}}}
%\newcommand{\bn}{{\mathfrak{n}} }
\newcommand{\sol}[1]{ \vskip 6pt{\color{blue} {\bf Solution:} #1} \vskip 6pt}

\newcommand{\bR}{{\mathbb{R}}}
\newcommand{\va}{{\vec{a}}}
\newcommand{\vb}{{\vec{b}}}
\newcommand{\vzero}{{\vec{0}}}
\newcommand{\vi}{{\vec{i}}}
\newcommand{\vj}{{\vec{j}}}
\newcommand{\vk}{{\vec{k}}}
\newcommand{\vh}{{\vec{h}}}
\newcommand{\vr}{{\vec{r}}}
\newcommand{\vu}{{\vec{u}}}
\newcommand{\vv}{{\vec{v}}}
\newcommand{\vw}{{\vec{w}}}
\newcommand{\vx}{{\vec{x}}}
\newcommand{\vy}{{\vec{y}}}
\newcommand{\vz}{{\vec{z}}}
\newcommand{\vT}{{\vec{T}}}
\newcommand{\vN}{{\vec{N}}}
\newcommand{\vF}{{\vec{F}}}
\newcommand{\vG}{{\vec{G}}}
\newcommand{\bF}{\mathbb{F}}
\newcommand{\bQ}{\mathbb{Q}}
\newcommand{\bZ}{\mathbb{Z}}
\newcommand{\bA}{\mathbb{A}}
\newcommand{\codim}{{\mathrm{codim}}}


\begin{document}
\title{Worksheet \#3 -- Math 6130\\Fall 2018}
\author{Due Monday, September 17th}

\maketitle
You may work in groups of up to 3.  Only one worksheet needs to be turned in per group.
\vskip 9pt
\noindent
{\bf 1.}  Consider $U = \bA^2 \setminus \{ 0 \}$.  Compute $R = \O_{\bA^2}(U)$ as a subring of $k(x,y)$, the fraction field of $\O_{\bA^2}(\bA^2) = k[x,y]$. 
\vskip 3pt
\emph{Hint:} There are different ways to do this.  One way is to use the fact that $k[x,y]$ is a UFD and suppose that $f/g \in R$, with $f, g \in k[x,y]$ having no common factors.  You should show that this expression has to be essentially unique.  On the other hand, we know $\dim Z(g) = 1$. 
\vskip 200pt
\noindent
{\bf 2.}  Prove that $U = \bA^2 \setminus \{ 0 \}$ is not isomorphic to an affine variety.  
\newpage
\noindent
{\bf 3.}  Consider the regular map $\phi : \bA^2 \to \bA^2$ defined by $\phi(x,y) = (x, xy)$.  Show that $\phi$ is dominant, birational, not an isomorphism and not finite.  
\vskip 250pt
\noindent
{\bf 4.}  Suppose that $X \subseteq \bA^n$ is a non-empty closed irreducible set such that $I(X) = (f_1, \dots, f_r)$ is generated by $r$ elements.  Show that $\codim_{\bA^n}(X) \leq r$.
\vskip 3pt
\emph{Hint:} Remember that $\codim_{\bA^n}(X) := \dim \bA^n - \dim X$.  We also may use that if $f \in k[Y]$ is a nonzero divisor on some algebraic variety $Y \subseteq \bA^n$, then $\dim Z_Y(f) = \dim Y - 1$ (Theorem 2.7.1 in the text).  The issue that makes this slightly tricky is that $Z_Y(f)$ is not necessarily itself an algebraic variety (there is no reason it must be irreducible).
\end{document}

